\documentclass[11pt,a4paper]{article}
\usepackage[margin=2cm]{geometry}
\usepackage{amsmath,amssymb}
\usepackage{xcolor}
\usepackage{enumitem}
\usepackage{tcolorbox}
\tcbuselibrary{breakable}

\setlength{\parindent}{0pt}
\setlength{\parskip}{4pt}

\newcommand{\pred}[1]{\textsc{#1}}

\title{\vspace{-1.5cm}Heuristics for Planning --- Q1 Notes}
\author{}
\date{}

\begin{document}
\maketitle
\vspace{-1cm}

\section*{The key insight}

In ordinary search, heuristics have to be supplied by a human who understands the domain. In planning the state and actions are described \emph{declaratively} --- each operator has explicit preconditions, add-list, delete-list --- so heuristics can be \textbf{derived automatically} from the problem description. This is what makes planning heuristics distinctive.

A generic recipe:
\begin{quote}
  \textit{Drop part of the problem to get an easier (``relaxed'') problem whose optimal cost is a lower bound on the original. Use that cost as $h$.}
\end{quote}
Provided the relaxation is genuinely easier (preferably polynomial), $h$ is cheap and admissible.

\section*{1.\ \ Relaxation heuristics (the big ones)}

\begin{tcolorbox}[colback=gray!5,colframe=gray!50!black,title=Ignore preconditions]
Drop all preconditions: every action is always applicable. The optimal relaxed plan length is the minimum number of actions whose add-lists cover the goal --- a set-cover problem. Useful but loose.
\end{tcolorbox}

\begin{tcolorbox}[colback=gray!5,colframe=gray!50!black,title=Ignore delete lists (delete-relaxation)]
Drop every action's delete-list, so nothing ever becomes false. The relaxed problem is monotone and solvable in polynomial time. The relaxed plan length is $h^+$.
\\[2pt]
\textbf{Why it works:} once a goal literal is achieved it stays achieved, so the planner never has to backtrack. \textbf{FastForward (FF)} uses exactly this as its heuristic --- one of the most successful planning systems.
\end{tcolorbox}

\begin{tcolorbox}[colback=gray!5,colframe=gray!50!black,title=Subgoal independence]
Plan for each conjunct of the goal in isolation. Then combine:
\begin{itemize}[leftmargin=*,itemsep=1pt]
  \item $h_{\max}(s) = \max_{g \in \text{goal}} \text{cost}(g)$ \ --- admissible, conservative.
  \item $h_{\text{add}}(s) = \sum_{g \in \text{goal}} \text{cost}(g)$ \ --- not admissible (double-counts), but usually more informative.
\end{itemize}
Assumes subgoals don't interact --- a useful lie.
\end{tcolorbox}

\section*{2.\ \ Use the structure of the goal}

\textbf{Landmarks.} A landmark is a fact (or action) that must hold in \emph{every} solution. Found by analysis of the goal and operators. They split planning into ``before-landmark'' and ``after-landmark'' phases and give a lower bound = number of landmarks not yet achieved.

\textbf{Helpful actions.} At each search node, only consider actions that appear in the relaxed plan from that node. Trades completeness for branching factor --- usually a big win.

\textbf{Means--ends analysis.} Pick an unmet goal literal $g$; pick an action whose add-list contains $g$; recurse on its preconditions. Old idea (GPS), still informs how heuristics rank actions.

\section*{3.\ \ Shrink the search space}

\textbf{Repeated-state detection.} Planning states are explicit sets of fluents --- equality is cheap to test. A closed-list / transposition table avoids exploring the same state twice. Big win for breadth-first style planning.

\textbf{Partial-order planning.} Only commit to an ordering between actions when there's a causal reason. Avoids exploring the $n!$ permutations of $n$ independent actions.

\textbf{State abstraction / pattern databases.} Project the state onto a subset of fluents; pre-compute optimal cost in that abstract space; use as $h$. Admissible by construction.

\section*{4.\ \ Decompose the problem}

\textbf{Hierarchical / HTN planning.} Plan with abstract operators first, then refine. Cuts the effective depth: planning a 1000-step low-level plan becomes planning a 10-step high-level plan, each step refined separately.

\textbf{Goal serialisation.} If two subgoals can be ordered such that solving the first never breaks the second (an ``ordered'' goal set), solve them one at a time. Finding such an ordering is itself a heuristic.

\section*{5.\ \ Reuse what you've already solved}

\textbf{Plan reuse / case-based planning.} Store solved plans; on a new problem, retrieve a similar past plan and adapt. Costly to maintain but pays off in repeated domains.

\textbf{Macro-operators.} Compile frequently used action sequences into single operators (e.g.\ ``open-and-walk-through-door''). Reduces tree depth at the cost of branching factor.

\section*{What to write in the exam answer}

A strong answer mentions \textbf{at least one heuristic from each of three categories}:

\begin{enumerate}[leftmargin=*,itemsep=2pt]
  \item A \textbf{relaxation} heuristic (ignore preconditions / ignore deletes / subgoal independence) --- and note that the planner can derive these \emph{automatically} from the operator descriptions; that's the whole point.
  \item A \textbf{structural} heuristic (landmarks, helpful actions, means--ends) --- which exploits the goal's syntactic structure.
  \item A \textbf{search-space} or \textbf{decomposition} idea (partial-order planning, HTN decomposition, repeated-state detection) --- which shrinks the tree rather than guiding it.
\end{enumerate}

End by linking back to the question: planning's edge over generic search is that operator descriptions expose enough problem structure for the planner to \textbf{compute its own heuristic}, rather than relying on a hand-written one.

\end{document}
